\section{折叠核 \texorpdfstring{$\Fold_m$}{Fold\_m} 及其 tail 分解}

\subsection{定义}

对 $m\ge 1$，定义
$$
\Fold_m:[0,2^m-1]\to X_m,\qquad \Fold_m(N)=(c_1(N),\dots,c_m(N)).
$$

\subsection{满射性}

对任意 $w\in X_m$，将其扩展为无限 Zeckendorf 合法序列 $(w_1,\dots,w_m,0,0,\dots)$，令
$$
N=V_m(w).
$$
则由唯一性可得 $\Fold_m(N)=w$，因此 $\Fold_m$ 满射。

\subsection{tail 计数表达式（一般上界 \texorpdfstring{$B$}{B}）}

先对一般上界 $B\in\NN$ 定义截断折叠映射
$$
\Fold_{m,B}:[0,B]\to X_m,\qquad \Fold_{m,B}(N)=(c_1(N),\dots,c_m(N)).
$$
并记 $s_{m,B}(w)=\abs{\Fold_{m,B}^{-1}(w)}$。对任意 $w\in X_m$，考虑所有满足 $\Fold_{m,B}(N)=w$ 的 $0\le N\le B$。写
$$
N=V_m(w)+T,\qquad
T=\sum_{k\ge m+1} c_k(N)\,F_{k+1}.
$$
由于 Zeckendorf 禁邻约束，必须满足边界相容条件
$$
w_m\,c_{m+1}(N)=0,
$$
且 tail 数码本身无相邻 $11$。因此
$$
s_{m,B}(w)=\abs{\Fold_{m,B}^{-1}(w)}
:=\#\Big\{T\in\mathcal{T}_{w_m}:\ T\le B-V_m(w)\Big\},
$$
其中 tail 和集合 $\mathcal{T}_0,\mathcal{T}_1$ 定义如下。

\begin{definition}[tail 和集合]\label{def:tail-sets}
对给定 $m\ge 1$，定义两类 admissible tail 数码集合
$$
\mathcal{D}_0=\Big\{(d_{m+1},d_{m+2},\dots)\in\{0,1\}^{\NN}:\ d_kd_{k+1}=0\ \forall k\ge m+1\Big\},
$$
$$
\mathcal{D}_1=\Big\{(d_{m+1},d_{m+2},\dots)\in\{0,1\}^{\NN}:\ d_{m+1}=0,\ d_kd_{k+1}=0\ \forall k\ge m+1\Big\}.
$$
对应的 tail 和集合定义为
$$
\mathcal{T}_\varepsilon=\Big\{\sum_{k\ge m+1} d_kF_{k+1}:\ (d_{m+1},d_{m+2},\dots)\in\mathcal{D}_\varepsilon\Big\},\qquad \varepsilon\in\{0,1\}.
$$
\end{definition}

该式是本文所有纤维结构的出发点：\textbf{纤维大小就是“反向枚举 tail”满足不等式约束的计数。}

\medskip
\noindent
本文的主对象取 dyadic 上界 $B_m=2^m-1$，并将
$$
\Fold_m:=\Fold_{m,B_m},\qquad s_m:=s_{m,B_m}
$$
作为简记。

